\documentclass[border=4pt]{standalone}
% Диаграмма вариантов использования системы автоматизации
% приёмной комиссии БНТУ. Чёрно-белое оформление по конвенции БНТУ.
%
% Компоновка:
%   * Акторы в одной вертикальной колонке (x=0). Каждый актор - узел только
%     для фигуры (без подписи в bbox), подпись отдельным узлом ниже.
%     Это позволяет стрелкам наследования упираться в фигуру, а не в текст.
%   * Каждый use case стоит на одной y-координате со своим основным актором
%     => стрелка ассоциации строго горизонтальная.
%   * Иерархия наследования с ромбовидным двойным родителем (Админ данных)
%     рисуется тонкими ортогональными огибалками слева от колонки акторов.
%
% Сборка: xelatex use-cases.tex (см. build-diagrams.sh).

\usepackage{fontspec}
\setmainfont{Times New Roman}

\usepackage{tikz}
\usetikzlibrary{positioning,calc,arrows.meta,fit,shapes.geometric}

% Размещение актора. Использование: \placeactor{name}{x}{y}{label}.
% Создаёт два узла: `name-fig` (только фигура) и `name-lbl` (подпись под
% фигурой). Стрелки наследования крепятся к name-lbl.south (старт ниже
% текста, чтобы линия не проходила через подпись) и name-fig.north.
% Стрелки ассоциаций - к name-fig.east.
% outer sep=5pt у обоих узлов: при использовании их анкоров (.south, .north,
% .east) TikZ возвращает точку на 5pt снаружи фактического содержимого,
% что автоматически даёт зазор между концами стрелок и текстом/фигурой.
\newcommand{\placeactor}[4]{%
  \begin{scope}[shift={(#2,#3)}]
    \node[minimum width=0.5cm, minimum height=1.0cm,
          inner sep=0pt, outer sep=5pt] (#1-fig) at (0,0) {};
    \draw[line width=0.5pt]
      (0, 0.40) circle (0.13)
      (0, 0.27) -- (0,-0.10)
      (-0.22, 0.15) -- (0.22, 0.15)
      (0,-0.10) -- (-0.18,-0.40)
      (0,-0.10) -- ( 0.18,-0.40);
    \node[anchor=north, font=\small, align=center,
          text width=3.2cm, inner sep=1pt, outer sep=5pt] (#1-lbl) at (0,-0.55) {#4};
  \end{scope}
}

\tikzset{
  usecase/.style={
    ellipse, draw, align=center, font=\small,
    minimum width=4.2cm, minimum height=1.3cm,
    inner sep=1pt, line width=0.4pt
  },
  inherit/.style={
    draw, -{Triangle[open, length=3mm, width=3mm]},
    line width=0.5pt
  },
  link/.style={
    draw, -{Latex[length=2mm]}, line width=0.4pt
  }
}

\begin{document}
\begin{tikzpicture}

% ============ Use cases (правая колонка, x=10) ============
% Порядок сверху вниз соответствует порядку акторов слева.
% Шаг 2 см, чтобы 3-строчные подписи акторов не накладывались на соседей.
\node[usecase] (uc1) at (10,12.0) {Управление\\учётными записями\\пользователей};
\node[usecase] (uc2) at (10, 9.0) {Настройка правил\\приёма и запуск\\пересчёта};
\node[usecase] (uc3) at (10, 6.0) {Валидация данных\\и разрешение\\удалений};
\node[usecase] (uc4) at (10, 3.0) {Ведение данных\\абитуриентов и заявок};
\node[usecase] (uc5) at (10, 1.2) {Аутентификация};
\node[usecase] (uc8) at (10,-0.6) {Просмотр данных\\всех сущностей};
\node[usecase] (uc6) at (10,-2.4) {Просмотр\\результатов отбора\\и выгрузка в Excel};
\node[usecase] (uc7) at (10,-4.2) {Просмотр журналов\\действий и входов};

% Граница системы. Title под верхней границей с fill=white, чтобы линия
% границы не пересекала текст подписи.
\node[draw, fit=(uc1)(uc2)(uc3)(uc4)(uc5)(uc8)(uc6)(uc7),
      inner xsep=12mm, inner ysep=12mm, line width=0.5pt] (system) {};
\node[anchor=north, font=\small\bfseries, fill=white, inner sep=2pt]
  at ($(system.north)+(0,-1mm)$) {Система автоматизации приёмной комиссии};

% ============ Акторы (левая колонка, x=0) ============
\placeactor{sa}{0}{12.0}{Суперадминистратор}
\placeactor{da}{0}{ 9.0}{Администратор данных}
\placeactor{au}{0}{ 6.0}{Аудитор}
\placeactor{op}{0}{ 3.0}{Оператор приёмной комиссии}
\placeactor{dv}{0}{-0.6}{Просматривающий}

% ============ Иерархия наследования (открытый треугольник к РОДИТЕЛЮ) ============
% Линейная цепочка: права каждой следующей роли - строгое надмножество прав
% предыдущей. SA = DA + управление пользователями, DA = AU + структура,
% AU = OP + аудит, OP = DV + ведение абитуриентов.
% Стрелки идут СВЕРХУ ВНИЗ от подписи верхнего актора (-lbl.south,
% чтобы линия не проходила через текст подписи) к фигуре нижнего
% актора (-fig.north). Треугольник у роли-родителя.
\draw[inherit] (sa-lbl.south) -- (da-fig.north);   % SuperAdmin         -> DataAdministrator
\draw[inherit] (da-lbl.south) -- (au-fig.north);   % DataAdministrator  -> Auditor
\draw[inherit] (au-lbl.south) -- (op-fig.north);   % Auditor            -> AdmissionsOperator
\draw[inherit] (op-lbl.south) -- (dv-fig.north);   % AdmissionsOperator -> DataViewer

% ============ Ассоциации актор --> вариант использования ============
% Все горизонтальные, без пересечений.
\draw[link] (sa-fig.east) -- (uc1.west);   % Суперадм -> Управление пользователями
\draw[link] (da-fig.east) -- (uc2.west);   % Админ    -> Настройка правил
\draw[link] (au-fig.east) -- (uc3.west);   % Аудитор  -> Валидация
\draw[link] (op-fig.east) -- (uc4.west);   % Оператор -> Ведение данных
\draw[link] (dv-fig.east) -- (uc5.west);   % Просматр -> Аутентификация
\draw[link] (dv-fig.east) -- (uc8.west);   % Просматр -> Просмотр всех данных
\draw[link] (dv-fig.east) -- (uc6.west);   % Просматр -> Просмотр результатов
\draw[link] (dv-fig.east) -- (uc7.west);   % Просматр -> Просмотр журналов

\end{tikzpicture}
\end{document}
